%==============================================================
%  lecons/algebre/euler_exp.tex
%  Définition limite de l'exponentielle matricielle
%==============================================================

\lecontitle{%
  $e^X = \lim_{k\to\infty}\!\left(I+\dfrac{X}{k}\right)^{\!k}$%
}{Algèbre linéaire — Définition limite et méthode d'Euler}

%=============================================================
%  § 1 — ÉNONCÉ
%=============================================================
\secnum{navyblue}{1}{Énoncé}

\begin{tcolorbox}[thmbox]
  \textbf{Théorème.}\enspace%
  \index{exponentielle de matrice!définition limite}
  Pour toute matrice $X \in \mathcal{M}_n(\mathbb{K})$, munie d'une norme
  sous-multiplicative $\|\cdot\|$~:
  \[
    e^X \;=\; \lim_{k \to +\infty} \left(I + \frac{X}{k}\right)^{\!k}.
  \]
\end{tcolorbox}

\rubriq{Rappel scalaire}
Pour $x \in \mathbb{R}$ ou $\mathbb{C}$, le résultat classique
$e^x = \lim_{k\to\infty}(1+x/k)^k$ s'obtient en posant
$f_k(x) = (1+x/k)^k$ et en remarquant que
$\ln f_k(x) = k\ln(1+x/k) \to x$.
La preuve matricielle ne peut utiliser le logarithme directement~;
elle passe par une comparaison terme à terme avec la série $e^X$.

\decsep{}

%=============================================================
%  § 2 — PREUVE
%=============================================================
\secnum{navyblue}{2}{Preuve}

\rubriq{Développement binomial}

\begin{mdframed}[style=proofbar]
  Posons $M = \|X\|$. Par le binôme de Newton (les puissances de $I$ et $X/k$
  commutent trivialement)~:
  \[
    \left(I + \frac{X}{k}\right)^{\!k}
    = \sum_{j=0}^{k} \binom{k}{j}\frac{X^j}{k^j}
    = \sum_{j=0}^{k} c_{k,j}\,X^j,
  \]
  où l'on pose
  \[
    c_{k,j}
    = \frac{k(k-1)\cdots(k-j+1)}{j!\,k^j}
    = \frac{1}{j!}\prod_{l=0}^{j-1}\!\left(1-\frac{l}{k}\right).
  \]

  \textit{Propriétés de $c_{k,j}$ (pour $k \geq j$)~:}
  \begin{itemize}
    \item $0 \leq c_{k,j} \leq \dfrac{1}{j!}$ \quad (le produit est $\leq 1$).
    \item $c_{k,j} \xrightarrow[k\to\infty]{} \dfrac{1}{j!}$ pour $j$ fixé.
  \end{itemize}
\end{mdframed}

\rubriq{Estimation de l'erreur}

\begin{mdframed}[style=proofbar]
  Par inégalité triangulaire~:
  \[
    \left\|\left(I+\frac{X}{k}\right)^{\!k} - e^X\right\|
    \;\leq\;
    \underbrace{\sum_{j=0}^{k}\!\left|c_{k,j}-\frac{1}{j!}\right|M^j}_{=:\,A_k}
    \;+\;
    \underbrace{\sum_{j=k+1}^{\infty}\frac{M^j}{j!}}_{=:\,B_k}.
  \]

  \textbf{Terme $B_k$.}\enspace
  C'est le reste de la série $e^M$ après le rang $k$~; elle converge
  absolument, donc $B_k \to 0$.

  \textbf{Terme $A_k$.}\enspace
  Soit $\varepsilon > 0$. Choisir $N$ tel que $2\sum_{j>N}M^j/j! < \varepsilon/2$.
  \begin{itemize}
    \item Pour $j > N$~: $|c_{k,j}-1/j!| \leq 2/j!$, donc
          $\sum_{j=N+1}^{k}|c_{k,j}-1/j!|M^j \leq 2\sum_{j>N}M^j/j! < \varepsilon/2$.
    \item Pour $j \leq N$~: en utilisant $1-\prod_l(1-l/k)\leq\sum_l l/k$,
          \[
            \left|c_{k,j}-\frac{1}{j!}\right|
            \leq \frac{1}{j!}\sum_{l=0}^{j-1}\frac{l}{k}
            \leq \frac{j^2}{2k\,j!},
          \]
          d'où $\sum_{j=0}^{N}|c_{k,j}-1/j!|M^j
          \leq \frac{1}{2k}\sum_{j=0}^N\frac{j^2 M^j}{j!}
          \xrightarrow[k\to\infty]{} 0$.
  \end{itemize}
  Pour $k$ assez grand, $A_k < \varepsilon$.
  Donc $A_k + B_k \to 0$. \hfill$\blacksquare$
\end{mdframed}

\decsep{}

%=============================================================
%  § 3 — PREUVE HEURISTIQUE VIA LE LOGARITHME
%=============================================================
\secnum{navyblue}{3}{Preuve heuristique via le logarithme matriciel}

\begin{mdframed}[style=proofbar]
  Pour $k$ assez grand, $\|X/k\| < 1$, et le logarithme matriciel est défini~:
  \[
    \log\!\left(I + \frac{X}{k}\right)
    = \sum_{m=1}^{\infty}\frac{(-1)^{m+1}}{m}\left(\frac{X}{k}\right)^{\!m}
    = \frac{X}{k} - \frac{X^2}{2k^2} + \cdots
    = \frac{X}{k} + O\!\left(\frac{1}{k^2}\right).
  \]
  Donc~:
  \[
    k\log\!\left(I+\frac{X}{k}\right) = X + O\!\left(\frac{1}{k}\right)
    \;\xrightarrow[k\to\infty]{}\; X.
  \]
  Comme $\exp$ est continue~:
  \[
    \left(I + \frac{X}{k}\right)^{\!k}
    = \exp\!\left(k\log\!\left(I+\frac{X}{k}\right)\right)
    \;\xrightarrow[k\to\infty]{}\; \exp(X). \hfill\square
  \]
  (Cette preuve est rigoureuse pour les matrices à condition de justifier
  que $\exp \circ (k\log(I+X/k)) \to \exp(X)$, ce que fait le §\,2.)
\end{mdframed}

\decsep{}

%=============================================================
%  § 4 — INTERPRÉTATIONS ET COROLLAIRES
%=============================================================
\secnum{navyblue}{4}{Interprétations et corollaires}

\begin{tcolorbox}[thmbox]
  \textbf{Corollaire 1 — Méthode d'Euler.}\enspace%
  \index{méthode d'Euler}
  Soit $y' = Xy$, $y(0) = y_0$. La méthode d'Euler avec $k$ pas de longueur
  $t/k$ donne~:
  \[
    y_k = \left(I + \frac{t}{k}X\right)^{\!k} y_0
    \;\xrightarrow[k\to\infty]{}\;
    e^{tX}\,y_0,
  \]
  la solution exacte.

  \medskip
  \textbf{Corollaire 2 — Formule de Trotter--Kato.}\enspace%
  \index{formule de Trotter}
  Pour $A, B \in \mathcal{M}_n(\mathbb{K})$ (non nécessairement commutatifs)~:
  \[
    e^{A+B} = \lim_{k\to\infty}\!\left(e^{A/k}\,e^{B/k}\right)^k.
  \]

  \medskip
  \textbf{Corollaire 3 — Continuité de $\exp$.}\enspace
  La formule $(I+X/k)^k \to e^X$ avec des estimées explicites fournit~:
  \[
    \|e^X - e^Y\| \leq \|X - Y\|\,e^{\max(\|X\|,\|Y\|)}.
  \]
\end{tcolorbox}

\rubriq{Preuve du corollaire 2 (formule de Trotter)}

\begin{mdframed}[style=proofbar]
  On utilise l'identité $e^{A/k}e^{B/k} = e^{(A+B)/k}e^{R_k}$ où
  $R_k = \log(e^{A/k}e^{B/k}) - (A+B)/k$. Le développement de Baker–Campbell–Hausdorff
  \index{Baker-Campbell-Hausdorff} donne $\|R_k\| = O(1/k^2)$, donc
  $k\|R_k\| \to 0$. Ainsi~:
  \[
    \left(e^{A/k}e^{B/k}\right)^k
    = e^{A+B}\,e^{kR_k + O(1/k)}
    \xrightarrow[k\to\infty]{} e^{A+B}.\hfill\blacksquare
  \]
\end{mdframed}

\rubriq{Interprétation~: intérêts composés}

\begin{mdframed}[style=exbar]
  Pour $X = rI_n$ ($r \in \mathbb{R}$)~:
  $\left(I + \frac{r}{k}I\right)^k = \left(1+\frac{r}{k}\right)^k I \to e^r I = e^{rI}$.

  Le cas scalaire $\left(1+\frac{r}{k}\right)^k \to e^r$ traduit la limite
  des intérêts composés~: capitaliser $r/k$ par période sur $k$ périodes
  tend, lorsque $k\to\infty$, vers la capitalisation continue $e^r$.
  La formule matricielle en est la généralisation naturelle.
\end{mdframed}

\decsep{}

%=============================================================
%  § 5 — PIÈGES ET EXERCICES
%=============================================================
\secnum{exgreen}{5}{Pièges et exercices}

\noindent{\bfseries\color{counterred} Pièges.}

\begin{mdframed}[style=cexbar]
  \textbf{Le binôme ne s'applique qu'à $I$ et $X/k$, pas à deux matrices
  quelconques.}\enspace
  Si $A \neq B$, $(A+B)^k \neq \sum_j\binom{k}{j}A^j B^{k-j}$ en général.
  Ici l'écriture $(I + X/k)^k$ fonctionne car $I$ et $X$ commutent toujours.

  \smallskip\noindent
  \textbf{La convergence est uniforme sur les bornés, pas sur tout
  $\mathcal{M}_n$.}\enspace
  Pour $\|X\| \leq R$ fixé, la vitesse de convergence est en
  $O(R^2/k)$~; mais pour $X$ quelconque, $k$ doit croître avec $\|X\|$.

  \smallskip\noindent
  \textbf{La formule de Trotter nécessite une preuve supplémentaire.}\enspace
  Contrairement à $(I+X/k)^k \to e^X$, la formule
  $(e^{A/k}e^{B/k})^k \to e^{A+B}$ n'est pas immédiate quand $AB \neq BA$~;
  elle repose sur Baker–Campbell–Hausdorff.
\end{mdframed}

\vspace{6pt}
\noindent{\bfseries\color{goldaccent} Exercices.}

\begin{mdframed}[style=exobar]
  \begin{enumerate}
    \item Retrouver $e^{\lambda I_n} = e^\lambda I_n$ à partir de la formule
          limite, puis calculer $\lim_{k\to\infty}(I + A/k)^k$ pour
          $A = \begin{pmatrix}0&1\\0&0\end{pmatrix}$.

    \item Montrer que $c_{k,j} \leq 1/j!$ et que $c_{k,j} \geq (1-j^2/(2k))/j!$
          pour $j \leq \sqrt{k}$. En déduire une vitesse de convergence~:
          \[
            \left\|\left(I+\frac{X}{k}\right)^{\!k} - e^X\right\|
            = O\!\left(\frac{\|X\|^2 e^{\|X\|}}{k}\right).
          \]

    \item En utilisant $(I+X/k)^k \to e^X$, prouver que
          $\det(e^X) = e^{\mathrm{tr}(X)}$ via
          $\det\!\bigl((I+X/k)^k\bigr) = \bigl(1 + \mathrm{tr}(X)/k
          + O(1/k^2)\bigr)^k$.

    \item Appliquer la méthode d'Euler avec $k=4$ pas au système
          $y'=Ay$, $y(0)=(1,0)^T$, $A=\begin{pmatrix}0&1\\-1&0\end{pmatrix}$,
          sur $[0,\pi]$. Comparer la valeur approchée à $e^{\pi A}(-I)^{1/2}$.

    \item Montrer que pour $A$ et $B$ qui commutent, la formule de Trotter
          se réduit immédiatement à $(I+X/k)^k \to e^X$ avec $X = A+B$.
          Expliquer pourquoi la preuve échoue sans commutativité.
  \end{enumerate}
\end{mdframed}

\leconfooter{%
  L.\ Euler (intérêts composés, 1748) —
  H.\ F.\ Trotter (formule du produit, 1959)%
}
